\documentclass[10pt]{article}
\usepackage[preprint]{tmlr}
\usepackage{amsmath,amssymb,booktabs,graphicx}
\usepackage[hidelinks]{hyperref}
\usepackage{url}
\hypersetup{pdftitle={What transfer rates do not tell you: conditioning practice in adversarial transfer evaluation}}
\title{What transfer rates do not tell you: conditioning practice in adversarial transfer evaluation}
\author{\name Anonymous Authors}
\def\month{09}
\def\year{2026}
\def\openreview{}
\setcounter{topnumber}{5}
\setcounter{bottomnumber}{5}
\setcounter{totalnumber}{10}
\renewcommand{\topfraction}{0.95}
\renewcommand{\bottomfraction}{0.95}
\renewcommand{\textfraction}{0.05}
\begin{document}
\maketitle
\begin{center}\small Working revision, 23 September 2026. Survey validation incomplete.\end{center}
\begin{abstract}
A provisional coding sheet for a preregistered, citation-ranked corpus of 80 papers marks 7 (8.75\%) as explicitly excluding clean target errors from the transfer denominator. This percentage is not yet a validated survey finding: a reliability audit invalidated the original repeatability calculation, and all judgments require evidence verification. We develop a reporting framework using the law-of-total-probability identity $\operatorname{PTR}_{st}=a_{st}\operatorname{CTR}_{st}+(1-a_{st})b_{st}$. Selected ImageNet reruns provide 27 internally verified summaries, but do not establish full replication of the corresponding papers. A single-paper extraction supplies 96 nominal lower bounds on newly induced error. A CIFAR-10 demonstration measures the identity's terms with exact counts. This working revision separates supported arithmetic from pending survey validation; it makes no claim about the prevalence of deficient evaluation in the literature.
\end{abstract}

\section{Introduction}
The provisional coding sheet marks 7 of 80 papers (8.75\%) as explicitly restricting the transfer denominator to clean-correct inputs. The other 73 are coded unclear, not no. These numbers describe an unvalidated sheet, not an established prevalence estimate. The central question remains useful independently of that percentage: what information is needed to reconstruct transfer conditional on source attack success?

Unconditional attacked error includes pre-existing target errors. Even after excluding those errors, pairwise transfer rate (PTR) and conditional transfer rate (CTR) concern different populations. Source ASR cannot generally connect them because its denominator is source-correct rather than jointly source-and-target-correct. We make these populations explicit and retain numerator and denominator counts.

Our contribution is a reporting framework and an empirical illustration of its consequences, not a new probability identity, attack, or discovery of transfer asymmetry. The identity follows directly from the law of total probability. Table~\ref{tab:conditioning} demonstrates a nonzero failed-source residual on CIFAR-10. The literature survey and its recomputability estimate remain provisional until a complete evidence audit and a separate blind recode are available.

\subsection{Related work and scope}
\label{sec:background}
Papernot et al.\ and Liu et al.\ established practical transfer-based black-box evaluations \citep{papernot,liu}. Directionality and optimization-dependent transfer are established phenomena \citep{demontis,dong}. Conditioning on both clean correctness and successful source attacks is not new: Inkawhich et al.\ explicitly use these conditions for their untargeted transfer evaluation \citep{inkawhich2019}. Our question concerns reporting sufficient to reconstruct CTR, never whether a named paper is wrong.

This framework complements general robustness-evaluation guidance \citep{carlini2019,tramer2020}, AutoAttack \citep{croce2020}, and RobustBench \citep{robustbench}; none of those benchmarks is claimed to have been run here. The mathematical setup concerns untargeted attacks against classifiers with compatible labels. Tangential corpus matches, including other tasks and surveys, need an explicit applicability ledger before substantive interpretation.

An audit of our own earlier repository implementation motivated the protocol. That implementation's label-space mismatch and treatment of pre-existing errors are not evidence about other papers, and its numerical results are excluded from this study.

\section{Why error is not attack success}
Let $f_t$ be the target predictor, $y$ the true label, and $A_s(x)$ an attack constructed using source $s$. Define clean error $e_t=\Pr[f_t(x)\ne y]$, newly induced error $q=\Pr[f_t(A_s(x))\ne y\mid f_t(x)=y]$, and persistence of existing error $r=\Pr[f_t(A_s(x))\ne y\mid f_t(x)\ne y]$. By partitioning on clean correctness,
\begin{equation}
\Pr[f_t(A_s(x))\ne y]=(1-e_t)q+e_t r.
\end{equation}
For the identity attack, $q=0$ and $r=1$, so adversarial error equals clean error. An illustrative classifier with 99\% clean error therefore has 99\% adversarial error under an identity attack and zero newly induced errors. This is an algebraic example, not a measurement from CIFAR-10.

The unconditional attacked-error rate and CTR condition on different populations, so their difference is not an estimate of inflation caused by clean errors. Equation~(1) attributes all-input attacked error to newly induced and persistent clean errors; PTR and CTR then add source-correctness and source-success conditions, with empty conditioning populations remaining undefined.

\subsection{Recorded estimands}
For a fixed evaluation set of $N$ examples, let $C_s(i)$ and $C_t(i)$ indicate clean correctness. Let $S_s(i)$ and $S_t(i)$ indicate misclassification of the source-generated attacked input. Define the pairwise eligible set $E_{st}=\{i:C_s(i)\land C_t(i)\}$ and successful-source subset $F_{st}=\{i\in E_{st}:S_s(i)\}$. We report
\begin{align}
\operatorname{ASR}_s &= \frac{\sum_{i:C_s(i)} S_s(i)}{\sum_i C_s(i)},\\
\operatorname{PTR}_{st} &= \frac{\sum_{i\in E_{st}} S_t(i)}{|E_{st}|},\\
\operatorname{CTR}_{st} &= \frac{\sum_{i\in F_{st}} S_t(i)}{|F_{st}|}.
\end{align}
PTR measures target failure on jointly clean-correct inputs whether or not the source attack succeeded. CTR measures transfer conditional on source success. To expose the arithmetic effect of conditioning, define pairwise source success $a_{st}=\Pr[S_s\mid E_{st}]$ and failed-source target error $b_{st}=\Pr[S_t\mid E_{st},\neg S_s]$. Partitioning the eligible set gives the exact identity
\begin{equation}
\operatorname{PTR}_{st}=a_{st}\operatorname{CTR}_{st}+(1-a_{st})b_{st}.
\end{equation}
Consequently, $\operatorname{CTR}_{st}=\operatorname{PTR}_{st}/a_{st}$ only when $b_{st}=0$. The conventional $\operatorname{ASR}_s$ above cannot be substituted for $a_{st}$ because its denominator includes source-correct inputs on which the target was already wrong. We report $a_{st}$ and $b_{st}$ for the conditioning analysis. The diagonal is a white-box control; its CTR is tautologically one whenever its denominator is positive.

The artifact tables carry numerator and denominator counts for every reported fraction. An empty denominator is undefined, not zero. Clean accuracy and attacked accuracy on the full evaluation set are also reported. Main cross-seed results use means and seed ranges; per-run artifact tables retain Wilson 95\% intervals under a binomial sampling interpretation for fixed checkpoints. Those intervals are descriptive for the finite CIFAR-10 test set and do not measure variability across training runs. Pairwise eligibility sets differ, so raw rates across model pairs are not strictly matched-population comparisons.

\paragraph{Boundary cases.}
The preceding rescaling statement assumes $0<a_{st}<1$. If $a_{st}=1$, $b_{st}$ is undefined but PTR equals CTR. If $a_{st}=0$, CTR is undefined. An empty eligible set leaves all pairwise rates undefined.
Likewise, $q$ is undefined when $e_t=1$ and $r$ when $e_t=0$; in these cases the clean-error decomposition is interpreted using joint probability masses, not numerical multiplication of an undefined conditional rate by zero.

\section{Survey methodology and coding scheme}
The corpus and six-field scheme were frozen at tag \texttt{survey-preregistration} before the recorded coding stage. Semantic Scholar was queried on 19 September 2026 for papers citing Liu et al.\ (arXiv:1611.02770) or Papernot et al.\ (arXiv:1602.02697). The filter retains metadata years 2018--2026 and titles or abstracts containing \emph{transfer} plus \emph{adversarial}, \emph{black-box}, or \emph{surrogate}. Eligible papers are ranked by citation count and the first 80 retained. The query, timestamp, raw paginated responses, and ranked list are preserved. Coverage of 2026 ends at retrieval; metadata years may differ from conference years. Citation ranking is purposive sampling, not a random sample of the literature.

The six yes/no/unclear fields concern separate source/target clean accuracy, clean-target-correct eligibility, source-success conditioning, absolute counts, a numerical $L_\infty$ or $L_2$ budget, and sufficient recomputation artifacts. Silence is unclear. Under the frozen rule, \emph{evaluated on correctly classified images} counts as yes for clean eligibility, although model-specific ambiguity must still be described. The frozen scheme is retained rather than retrospectively altered.

\paragraph{Reliability correction.}
The historical first and second sheets were generated using the same hard-coded judgment sets. Their agreement is a self-comparison, not a blind repeatability assessment. The previously reported $\kappa=1.00$ is withdrawn; no valid $\kappa$ is currently available. Waiving the planned seven-day delay does not waive blinding. A separate blind recode of the fixed 16-paper sample, followed by comparison against a frozen, evidence-audited first pass, remains necessary. If a defined per-field $\kappa$ is below 0.7, the scheme must be tightened and all 80 papers recoded; degenerate $\kappa$ is undefined, not perfect agreement.

\section{Survey results}
Table~\ref{tab:survey} preserves the historical sheet's descriptive counts for auditability. Evidence snippets generated by keyword search sometimes cite titles or general discussion rather than the relevant evaluation protocol. Consequently, all 480 judgments need verification against the paper and supplement, with field-specific page, section or table evidence. In particular, a failure to find a statement cannot support no.

\begin{table}[t]
\centering
\caption{Provisional, unvalidated coding counts ($N=80$). These are not confirmed prevalence estimates.}
\label{tab:survey}
\begin{tabular}{lrrr}
\toprule
Field & Yes & No & Unclear\\
\midrule
Separate clean accuracy & 5 & 0 & 75\\
Clean-correct eligibility & 7 & 0 & 73\\
Source-success conditioning & 1 & 1 & 78\\
Absolute counts & 0 & 57 & 23\\
Numerical norm budget & 46 & 9 & 25\\
Recomputation artifacts & 3 & 0 & 77\\
\bottomrule
\end{tabular}
\end{table}

The artifact audit examined five candidate repositories, not all 80 papers exhaustively. Thus 3/80 is a provisional positive-code fraction, not a validated recomputability rate; the remaining 77 are not demonstrated unavailable. No corpus-wide rate of reusable per-example outputs is established. A completed audit must report recomputability as a result regardless of its value. If most papers condition correctly, that is the result to publish. The frozen corpus must not be changed to obtain a stronger conclusion.

\section{Recomputation and analytic bounds}
\subsection{Selected artifact reruns}
Three selected ImageNet configurations yielded 27 source--target summaries, each from 1,000 images. These runs do not cover every experiment in the three papers. Two use PyTorch ports of TensorFlow~1 implementations; original-runtime numerical equivalence has not been established. Common evaluation images and converted target weights came through the SSA checkout. Shared architecture names do not establish identical checkpoint bytes or evaluation populations.

For SSA, source labels are zero-based while converted target predictions are one-based, so target correctness compares predictions to label plus one. The SSA Inception-v3 target differs from the source implementation and is not a white-box diagonal. The other two runs use one-based labels for both. A dependency-free verifier reconstructs all 27 archived count/rate summaries from their per-example predictions using these conventions.

Across these selected summaries PTR ranges from 14.64\% to 100\%, CTR from 14.68\% to 100\%, and $a_{st}$ from 99.66\% to 100\%. The failed-source stratum is empty in 18 summaries, leaving $b_{st}$ undefined; in the remaining nine it is zero. Twenty entries have nominal published comparators, with unconditional-rate differences of $-3.9$ to $+5.5$ percentage points (median $-0.65$). Without matched checkpoints, data, preprocessing, and original-runtime validation, those differences are not evidence of replication fidelity or a conditioning effect. Arithmetic verification alone does not validate perturbation budgets or attack implementations.

\subsection{Bounds without per-example predictions}
Write $u_t$ for unconditional attacked error and $e_t$ for target clean error on the same population. Equation~1 gives the sharp bounds
\begin{equation}
\max(0,u_t-e_t)\leq (1-e_t)q\leq\min(u_t,1-e_t).
\end{equation}
The quantity $e_t$ bounds the absolute probability mass attributable to persistent clean error, not its fraction of $u_t$. For $u_t>0$, that fraction is at most $\min(1,e_t/u_t)$. For $e_t<1$, dividing the lower bound by $1-e_t$ bounds target-clean-conditional newly induced error, not CTR. Recovering CTR additionally requires source-side information.

The current extraction covers 96 off-diagonal rates from Table~1 of Wu et al.\ \citep{wu2020}, not 96 papers or a corpus-wide distribution. Twenty-four of these rates use that paper's ensemble target, which averages all four models including the source, so only 72 are disjoint single-model transfers; the paper's Table~3 contains further eligible cells not yet included. Nominal lower bounds range from 1.2 to 77.9 percentage points (median 34.5). A separate conservative calculation assumes rounding to the nearest displayed unit, including an interval of $[99.5,100]$ for the clean entry printed as $100\%$. Its lower bounds range from 0.95 to 77.35 percentage points, with median 34.175 and mean 35.143. These intervals reflect assumed reporting precision, not sampling uncertainty. The rates are correlated within-paper measurements. An independent second extraction and an eligibility/population-alignment ledger for every paper remain necessary before presenting a corpus result. The current extraction is a worked case only.

\section{CIFAR-10 demonstration}
We use three independently trained checkpoints each of ResNet-18 (R), VGG16 (V), and MobileNetV2 (M), all with ten outputs and compatible class order. Checkpoint selection uses a fixed 45,000/5,000 training/validation split; evaluation uses all 10,000 test images. Mean clean accuracy is 95.11\%, 90.63\%, and 90.31\%, respectively. Normalization occurs inside the model, leaving budgets in raw-pixel coordinates. Full training, FGSM, matched-learning-rate, and convergence details appear in Appendix~\ref{app:cifar}.

The demonstration uses source-selected PGD with 40 steps, five restarts, and step size one quarter of the $L_\infty$ radius. The same attacked inputs are evaluated on every target; target predictions never select an iterate. Per-example predictions and counts support each reported denominator.

At $2/255$, Table~\ref{tab:conditioning} separates jointly clean-correct eligibility from source success. Five directions have mean $b_{st}\leq0.68\%$. VGG16-to-MobileNetV2 has mean $b_{st}=2.83\%$ and observed mean CTR 1.19 percentage points below mean PTR divided by mean $a_{st}$. These averages are descriptive: the identity is evaluated per seed, not by multiplying independently averaged terms. Conventional source ASR and PTR alone cannot reconstruct CTR. The experiment illustrates the distinction, not its frequency in other architectures or datasets.

\setcounter{table}{5}
\begin{table}[t]
\centering
\caption{$2/255$ PGD conditioning decomposition. Count tuples list seeds 0/1/2; rates are mean [range] percentages.}
\label{tab:conditioning}
\resizebox{\columnwidth}{!}{%
\begin{tabular}{lrr}
\toprule
Direction & Pairwise eligible $|E|$ & Source success $|F|$\\
\midrule
R$\rightarrow$V & 8905/8874/8923 & 8381/8318/8424\\
R$\rightarrow$M & 8831/8894/8873 & 8307/8339/8374\\
V$\rightarrow$R & 8905/8874/8923 & 6457/6257/6303\\
V$\rightarrow$M & 8583/8589/8620 & 6139/5976/6003\\
M$\rightarrow$R & 8831/8894/8873 & 8822/8879/8857\\
M$\rightarrow$V & 8583/8589/8620 & 8574/8574/8604\\
\bottomrule
\end{tabular}
}
\vspace{3pt}
\resizebox{\columnwidth}{!}{%
\begin{tabular}{lrrrr}
\toprule
Direction & Pair ASR & Fail-target & PTR & CTR\\
\midrule
R$\rightarrow$V & 94.09 [93.73,94.41] & 0.00 [0.00,0.00] & 5.84 & 6.21\\
R$\rightarrow$M & 94.07 [93.76,94.38] & 0.25 [0.19,0.36] & 21.61 & 22.96\\
V$\rightarrow$R & 71.22 [70.51,72.51] & 0.68 [0.42,0.96] & 15.45 & 21.42\\
V$\rightarrow$M & 70.25 [69.58,71.53] & 2.83 [2.71,2.99] & 23.29 & 31.97\\
M$\rightarrow$R & 99.85 [99.82,99.90] & 0.00 [0.00,0.00] & 10.79 & 10.81\\
M$\rightarrow$V & 99.84 [99.81,99.90] & 0.00 [0.00,0.00] & 4.61 & 4.62\\
\bottomrule
\end{tabular}
}
\end{table}

\section{Limitations}
The survey is not publication-ready: its evidence audit and blind reliability check remain incomplete. Citation-ranked selection, partial-year coverage, and task heterogeneity limit generalization. The bounds extraction covers one paper, and the selected artifact reruns lack validated original-runtime equivalence. These limitations cannot be repaired by wording alone.

Code, manuscript, and reproducibility snapshot are archived under concept DOI \href{https://doi.org/10.5281/zenodo.22737839}{10.5281/zenodo.22737839}; the historical results were released as v1.2.0. This working revision corrects its unsupported survey and reliability claims and is not a new validated release. The per-example experiment records are deposited separately under concept DOI \href{https://doi.org/10.5281/zenodo.22756878}{10.5281/zenodo.22756878}. The final record contains nine checkpoints, full-test predictions, exact eligibility counts, perturbation norms, manifests, and checkpoint and source hashes. The three seed triads provide evidence about variability under the specified training procedure, but three draws support descriptive ranges rather than precise population-level inference. Per-run Wilson intervals describe fixed-checkpoint binomial variation and do not capture training variability. The same 10,000 test images recur across seeds, so image-level counts must not be pooled as 30,000 independent observations.

The CIFAR-10 demonstration covers three standard CNN families on CIFAR-10 and untargeted $L_\infty$ attacks. It does not establish behavior on larger datasets, transformers, adversarially trained models, targeted attacks, or physical settings. It also does not compare against every modern transfer attack. The VGG16 learning-rate adjustment was selected using validation results after the common recipe failed. A seed-0 matched-learning-rate control preserves its strongest asymmetries, but only one control seed was run and some individual entries change substantially; optimization and architecture remain entangled. Random-noise draws are independently generated for each nominal source row, so noise-control rows are reproducible but do not share one common perturbation realization. The retained manifests record a dirty working tree at execution time; they preserve prediction checksums, checkpoint hashes, the recorded commit, and hashes of the executed source files, while the clean release tag archives the final corrected code.

The original artifact is a motivating self-correction, not evidence that similar errors are common elsewhere. The present results do not imply that one architecture is intrinsically more transferable: architecture, parameterization, training dynamics, and clean accuracy vary together. PTR and CTR answer different questions, and neither should be selected after seeing which produces the more dramatic result.

\section{Conclusion}
Explicit conditioning populations and counts make transfer rates interpretable. The probability identities and the CIFAR-10 demonstration show why source ASR and PTR generally cannot reconstruct CTR. Selected ImageNet prediction archives pass arithmetic verification, but this does not establish full replication. Survey counts, recomputability prevalence, and corpus-wide bounds remain pending evidence validation; no valid reliability estimate is currently available. A completed study must report the validated result, including a result that most papers condition correctly.

\appendix
\section{CIFAR-10 protocol and supplementary results}
\label{app:cifar}
\subsection{Data and model training}
The pipeline uses the standard CIFAR-10 class order \citep{cifar}. A fixed split seed of 1729 partitions the 50,000 training images into 45,000 training and 5,000 validation images. Validation data receive no random augmentation. Training uses random crops with four-pixel padding and random horizontal flips. The 10,000 test images are reserved for final evaluation; checkpoint selection uses validation accuracy only.

All networks have ten output classes and are initialized without ImageNet weights. The ResNet-18 stem uses a $3\times3$ convolution at stride one and no initial max pool. VGG16 uses its convolutional feature stack, global $1\times1$ pooling, and a single 512-to-10 linear classifier. MobileNetV2 changes its initial convolution to stride one. These adaptations, especially the compact VGG classifier, must be stated when comparing with other implementations.

Inputs to the attack interface are floating-point pixels in $[0,1]$. Channel normalization with mean $(0.4914,0.4822,0.4465)$ and standard deviation $(0.2470,0.2435,0.2616)$ occurs inside the model, so gradients include preprocessing while perturbation budgets remain in raw-pixel coordinates. Checkpoints contain architecture, class order, split indices, training seed, selected epoch, and model state. Evaluation rejects incompatible checkpoint metadata.

Training uses 200 epochs, SGD with momentum 0.9, weight decay $5\times10^{-4}$, cosine learning-rate decay, and batch size 128. ResNet-18 and MobileNetV2 use initial learning rate 0.1. That rate caused the VGG16 loss to remain near $\log 10$ and validation accuracy near chance, so we stopped that failed run and selected 0.01 using validation performance only. We then froze 0.01 for all three VGG16 seeds before test evaluation. The independently trained seeds are 0, 1, and 2; the split seed remains 1729.

\subsection{Attacks and controls}
FGSM uses one signed cross-entropy gradient step with projection into the pixel range. PGD uses random initialization within the intersection of the pixel range and the $L_\infty$ ball, 40 steps, absolute step size $2/255$, and five restarts at the primary budget $8/255$. Across iterates and restarts, the implementation prioritizes source misclassification, then larger source loss. Target predictions never influence this selection. This source-selected construction may differ from final-iterate PGD in its transfer behavior, so the selection rule is part of the attack specification.

The clean control leaves inputs unchanged. The random-noise control samples independent uniform perturbations within the same $L_\infty$ budget and clips to the pixel range. The empirical scope is limited to the $L_\infty$ threat model. We do not mix CW-$L_2$ results with these measurements or imply that equal numeric radii across norms represent equal threat models.

For every source and attack, the same generated inputs are evaluated on every target, including the source. Models remain in evaluation mode. Runtime checks reject nonfinite pixels, invalid pixel ranges, and norm violations. A saved example records its test index, true label, predictions, and actual perturbation. Visualization does not presume attack success or establish perceptual invisibility.

\subsection{Reproducibility and sensitivity analysis}
A seeded permutation specifies test indices, and the identical ordered sample list was used for every pair. All reported final evaluations used the full 10,000-image CIFAR-10 test set; the earlier 1,000-example runs were labeled as pilots and excluded from the empirical results. Each completed run saves arguments, software versions, commit identifier, checkpoint SHA-256 hashes, test indices, elapsed time, clean and attacked per-example predictions, and per-example $L_2$ and $L_\infty$ norms. Tables are generated from measured summaries rather than embedded constants.

We evaluate $L_\infty$ radii $2/255$, $4/255$, and $8/255$. The PGD step size is one quarter of the radius. At $8/255$, seed 0 uses 10, 40, and 100 steps with one restart and 100 steps with five restarts; seeds 1 and 2 repeat the 10-step, one-restart and 100-step, five-restart endpoints. At $2/255$, seed 0 also uses 100 steps and five restarts. Attack hyperparameters are not selected to maximize target transfer. We show training seeds individually and summarize seed-level estimates; we do not pool repeated test images as independent observations.

To test whether the architecture-specific VGG16 learning rate alone explains the largest asymmetries, we train seed-0 ResNet-18 and MobileNetV2 controls at 0.01 for the same 200 epochs and rerun the full-test $8/255$ protocol with the existing VGG16 seed-0 checkpoint. This is a prespecified sensitivity comparison; it is not pooled with the three-seed primary estimates.

\subsection{Clean accuracy and primary attack strength}
Table~\ref{tab:clean} reports validation-selected checkpoint performance on the untouched test set. Variation across the three training seeds is small relative to the differences among architectures. At the primary $8/255$ budget, 40-step, five-restart PGD attains 100\% source ASR for every ResNet-18 and MobileNetV2 checkpoint and at least 99.98\% for every VGG16 checkpoint. The diagonal therefore confirms that the primary attack is strong against its source models.

Uniform random noise at the same $8/255$ radius produces mean off-diagonal PTR between 1.45\% and 7.89\%, below the corresponding adversarial results. Noise CTR is based on much smaller source-success subsets and is not used to rank model pairs. Because each nominal source row uses a separately seeded noise realization, these controls establish a non-adversarial baseline within each row rather than a paired comparison among noise rows.

\begin{table}[t]
\centering
\caption{Clean test accuracy (\%) by training seed. The final column is the mean and sample standard deviation over three checkpoints.}
\label{tab:clean}
\resizebox{\columnwidth}{!}{%
\begin{tabular}{lrrrr}
\toprule
Model & Seed 0 & Seed 1 & Seed 2 & Mean $\pm$ SD\\
\midrule
ResNet-18 & 95.11 & 95.01 & 95.20 & 95.11 $\pm$ 0.10\\
VGG16 & 90.59 & 90.41 & 90.88 & 90.63 $\pm$ 0.24\\
MobileNetV2 & 89.97 & 90.57 & 90.38 & 90.31 $\pm$ 0.31\\
\bottomrule
\end{tabular}
}
\end{table}

\begin{table}[t]
\centering
\caption{Mean source ASR (\%) over three seeds.}
\label{tab:asr}
\resizebox{\columnwidth}{!}{%
\begin{tabular}{lrrrr}
\toprule
Source & FGSM $8/255$ & PGD $2/255$ & PGD $4/255$ & PGD $8/255$\\
\midrule
ResNet-18 & 58.54 & 94.47 & 99.91 & 100.00\\
VGG16 & 87.31 & 71.73 & 97.93 & 99.98\\
MobileNetV2 & 92.03 & 99.85 & 100.00 & 100.00\\
\bottomrule
\end{tabular}
}
\end{table}

\subsection{Directionality at the primary budget}
Table~\ref{tab:primary} shows off-diagonal PGD PTR at $8/255$. Transfer is strongly directional. VGG16 examples transfer to ResNet-18 at 92.73\% on average, whereas the reverse direction reaches 18.09\%. MobileNetV2-to-VGG16 is 25.94\%, whereas VGG16-to-MobileNetV2 is 95.23\%. The ordering persists across all three seeds, so it is not produced by one checkpoint draw. These are associations among the trained architectures and recipes; differences in capacity, optimization, and clean accuracy preclude an architecture-only causal interpretation.

\begin{table}[t]
\centering
\caption{PGD PTR (\%) at $8/255$, reported as mean [minimum, maximum] over three training seeds.}
\label{tab:primary}
\resizebox{\columnwidth}{!}{%
\begin{tabular}{lrrr}
\toprule
Source & ResNet-18 & VGG16 & MobileNetV2\\
\midrule
ResNet-18 & --- & 18.09 [16.73,19.38] & 72.27 [71.76,72.73]\\
VGG16 & 92.73 [92.40,92.99] & --- & 95.23 [94.80,95.61]\\
MobileNetV2 & 67.60 [63.46,70.61] & 25.94 [25.73,26.29] & ---\\
\bottomrule
\end{tabular}
}
\end{table}

\subsection{Budget and conditioning sensitivity}
Transfer falls sharply with the perturbation budget, but the directional ordering is broadly stable (Figure~\ref{fig:budget}). VGG16-to-ResNet-18 drops from 92.73\% at $8/255$ to 15.45\% at $2/255$; ResNet-18-to-MobileNetV2 drops from 72.27\% to 21.61\%. These changes rule out treating one reported transfer percentage as a property of a model pair independent of its threat model.

\begin{figure}[t]
\centering
\includegraphics[width=\columnwidth]{../../artifacts/study/transfer_summary.pdf}
\caption{Mean PGD PTR across three seeds versus $L_\infty$ budget. Error bars span the minimum and maximum seed values.}
\label{fig:budget}
\end{figure}

At $2/255$, VGG16 source ASR is only 71.05--72.98\% (Table~\ref{tab:asr}). Table~\ref{tab:conditioning} applies the exact conditioning identity with absolute counts. Conditioning is nearly mechanical in five directions: mean $b_{st}$ is at most 0.68\%, and mean CTR is within 0.27 points of the value obtained by dividing mean PTR by mean $a_{st}$. Across individual seeds, the largest corresponding deviation is 0.40 points. VGG16-to-MobileNetV2 is the exception, with mean $b_{st}=2.83\%$ and observed CTR 1.19 points below the zero-residual value. This exception confirms that the successful- and failed-source strata are not interchangeable. More fundamentally, conventional source ASR and PTR cannot determine CTR even when the residual is negligible, because conventional ASR uses a broader, target-independent denominator than $a_{st}$.


\section{Additional CIFAR-10 controls}
The stronger 100-step, five-restart $8/255$ endpoint changes off-diagonal PTR by at most 1.70 percentage points relative to the primary 40-step, five-restart setting across all seeds. At $2/255$ on seed 0, increasing from 40 to 100 steps at five restarts changes all off-diagonal PTR entries by at most 0.27 points and CTR entries by at most 0.40 points.

\begin{table}[!htbp]
\centering
\caption{Seed-0 PGD PTR (\%) at $8/255$ before and after matching the initial learning rate.}
\label{tab:lrcontrol}
\begin{tabular}{lrr}
\toprule
Direction & Primary & Matched LR\\
\midrule
ResNet-18 $\rightarrow$ VGG16 & 19.38 & 32.71\\
VGG16 $\rightarrow$ ResNet-18 & 92.99 & 95.98\\
MobileNetV2 $\rightarrow$ VGG16 & 25.73 & 18.55\\
VGG16 $\rightarrow$ MobileNetV2 & 94.80 & 93.12\\
\bottomrule
\end{tabular}
\end{table}

\begin{table}[!htbp]
\centering
\caption{FGSM PTR (\%) at $8/255$, mean [minimum, maximum] over three seeds.}
\label{tab:fgsm}
\resizebox{\columnwidth}{!}{%
\begin{tabular}{lrrr}
\toprule
Source & ResNet-18 & VGG16 & MobileNetV2\\
\midrule
ResNet-18 & --- & 27.80 [26.21,30.01] & 58.30 [56.73,59.43]\\
VGG16 & 55.27 [54.14,56.75] & --- & 64.88 [64.38,65.71]\\
MobileNetV2 & 47.37 [47.20,47.55] & 29.34 [28.09,30.55] & ---\\
\bottomrule
\end{tabular}}
\end{table}

\clearpage
\bibliography{main}
\bibliographystyle{tmlr}
\end{document}
